\chapter{Mã nguồn lõi: Tarjan SCC và BFS bán kính ảnh hưởng}
\label{app:dsa}

Hai giải thuật đồ thị của pipeline (kho \texttt{K8sSecDaP-pipeline},
\texttt{libdsa/src/graph/}) phục vụ phát hiện dịch chuyển ngang: Tarjan tìm thành phần
liên thông mạnh (chu trình giao tiếp bất thường), còn BFS đo bán kính ảnh hưởng (số pod
một nguồn tiếp cận được trong $k$ bước). Cơ sở lý thuyết đã trình bày ở
Chương~\ref{ch:foundations}; phần này trích phần lõi của hiện thực.

\section{Tarjan SCC (hiện thực lặp)}

Để tránh tràn ngăn xếp lời gọi trên đồ thị lớn, Tarjan được hiện thực bằng DFS
\emph{lặp} (dùng ngăn xếp tường minh) thay vì đệ quy; mỗi đỉnh giữ chỉ số khám phá
\texttt{disc} và liên kết thấp \texttt{low}, một SCC được kết xuất khi
\texttt{low[u] == disc[u]}.

\begin{lstlisting}[language=C++,caption={Điểm mấu chốt của Tarjan: cập nhật low-link và kết xuất SCC (\texttt{tarjan.cpp}).}]
// Khi quay lui tu dinh con v ve cha u:
low[cur.u] = std::min(low[cur.u], low[v]);
// Khi gap dinh v da tham va con tren stack (canh nguoc):
if (on_stack[v]) low[cur.u] = std::min(low[cur.u], disc[v]);
// u la goc cua mot SCC:
if (low[cur.u] == disc[cur.u]) {
    Component scc;
    uint32_t w;
    do { w = stk.top(); stk.pop(); on_stack[w] = false; scc.push_back(w); }
    while (w != cur.u);
    components.push_back(std::move(scc));
}
\end{lstlisting}

\section{BFS bán kính ảnh hưởng (k-hop)}

BFS duyệt theo lớp từ địa chỉ nguồn, dừng ở độ sâu \texttt{max\_hops}; tập đỉnh tới được
chính là \emph{bán kính ảnh hưởng} (blast radius) của nguồn~--~thước đo mức độ lan rộng
nếu nguồn đó tiếp tục dịch chuyển ngang.

\begin{lstlisting}[language=C++,caption={Duyệt BFS theo lớp, dừng ở $k$ bước (\texttt{bfs.cpp}).}]
ReachabilityResult BFSReachability::reachable_within(
    const DirectedGraph& g, uint32_t source, int max_hops) {
    std::queue<uint32_t> q;  q.push(source);
    std::unordered_map<uint32_t,int> dist; dist[source] = 0;
    while (!q.empty()) {
        uint32_t u = q.front(); q.pop();
        int d = dist[u];
        if (d >= max_hops) continue;            // khong mo rong qua k buoc
        for (uint32_t v : g.neighbors(u))
            if (dist.find(v) == dist.end()) { dist[v] = d + 1; q.push(v); }
    }
    return dist;   // tap dinh toi duoc + so buoc -> ban kinh anh huong
}
\end{lstlisting}
